\section*{Capítulo \ref{ch:g9:linfunc} --- Funções lineares e afins}

\begin{solution}{exo:g9:linfunc:1}
Lineares: $f$ e $k$ (da forma $ax$). Afins, mas não lineares: $g$
(\omterm{prop:g9:linfunc:affinegraph}{coeficiente angular} $3$, coeficiente linear $-5$) e $m$ (constante,
\omterm{prop:g9:linfunc:affinegraph}{coeficiente angular} $0$). Nem uma coisa nem outra: $h$, por causa do
quadrado.
\end{solution}

\begin{solution}{exo:g9:linfunc:2}
$f(0) = -3$; $f(4) = 5$; $f(-1) = -5$. E $f(x) = 9$ significa
$2x - 3 = 9$, então $2x = 12$ e $x = 6$.
\end{solution}

\begin{solution}{exo:g9:linfunc:3}
$a = \frac{15}{6} = 2.5$, então $f(x) = 2.5x$. Daí, $f(10) = 25$ e
$f(-2) = -5$.
\end{solution}

\begin{solution}{exo:g9:linfunc:4}
$f$ passa pela origem com \omterm{prop:g9:linfunc:affinegraph}{coeficiente angular} $2$; $g$ tem o mesmo
\omterm{prop:g9:linfunc:affinegraph}{coeficiente angular}, mas cruza o eixo vertical em $3$; $h$ desce a
partir de $(0, 4)$ com \omterm{prop:g9:linfunc:affinegraph}{coeficiente angular} $-1$. Os gráficos de $f$ e
de $g$ são \emph{\omterm{def:g6:lines:perp}{retas paralelas}} (mesmo \omterm{prop:g9:linfunc:affinegraph}{coeficiente angular},
coeficientes lineares diferentes).
\end{solution}

\begin{solution}{exo:g9:linfunc:5}
Para $f$: \omterm{prop:g9:linfunc:affinegraph}{coeficiente angular} $\frac{11 - 5}{3 - 1} = 3$, e depois
$5 = 3 \times 1 + b$ dá $b = 2$: $f(x) = 3x + 2$.

Para $g$: \omterm{prop:g9:linfunc:affinegraph}{coeficiente angular} $\frac{0 - 4}{2 - 0} = -2$, e $g(0) = 4$
já é o coeficiente linear: $g(x) = -2x + 4$.
\end{solution}

\begin{solution}{exo:g9:linfunc:6}
Aumento de $25\%$: $120 \times 1.25 = 150$.

Desconto de $40\%$: $65 \times 0.6 = 39$.

Preço original $p$ com $p \times 0.75 = 69$:
$p = \frac{69}{0.75} = 92$.
\end{solution}

\begin{solution}{exo:g9:linfunc:7}
\emph{1.} $f(x) = 0.05x + 10$.

\emph{2.} O plano B fica mais barato quando $25 < 0.05x + 10$, isto é,
$15 < 0.05x$, isto é, $x > 300$. A partir de $301$ minutos ($5$ horas)
por mês, o plano B ganha.
\end{solution}

\begin{solution}{exo:g9:linfunc:8}
\emph{1.} Depois de uma hora: $2000 \times 1.1 = 2200$. Depois de duas
horas: $2200 \times 1.1 = 2420$.

\emph{2.} O segundo aumento se aplica a $2200$, e não a $2000$: crescer
$10\%$ duas vezes multiplica por $1.1^2 = 1.21$, um aumento de $21\%$,
e não de $20\%$.
\end{solution}

\begin{solution}{exo:g9:linfunc:9}
\emph{1.} \omterm{prop:g9:linfunc:affinegraph}{Coeficiente angular}: $\frac{1 - 3}{6 - 2} = \frac{-2}{4} =
-0.5$: a função é decrescente.

\emph{2.} $f(x) = -0.5x + b$ com $f(2) = 3$: $-1 + b = 3$, então
$b = 4$ e $f(x) = -0.5x + 4$. Saída $0$: $-0.5x + 4 = 0$ dá $x = 8$.
\end{solution}

\begin{solution}{exo:g9:linfunc:10}
\emph{1.} As duas variações multiplicam o preço por
\[
\left(1 + \frac{t}{100}\right)\left(1 - \frac{t}{100}\right)
= 1 - \frac{t^2}{10000}
\]
(terceiro \omterm{def:g2:mult:def}{produto} notável com $a = 1$, $b = \frac{t}{100}$).

\emph{2.} Para $0 < t \leq 100$, $\frac{t^2}{10000} > 0$, então o fator
é menor que $1$: o preço final fica abaixo do original. Para $t = 10$:
fator $1 - \frac{100}{10000} = 0.99$, ou seja, uma queda total de
$1\%$.
\end{solution}

\begin{solution}{pb:g9:linfunc:1}
\textbf{1.} \omterm{prop:g9:linfunc:affinegraph}{Coeficiente angular}:
$\frac{f(100) - f(0)}{100 - 0} = \frac{212 - 32}{100} = 1.8$;
coeficiente linear $f(0) = 32$. Logo,
\[
f(C) = 1.8\,C + 32 .
\]

\textbf{2.} $f(20) = 68\,^\circ$F (o famoso dia ameno);
$f(37) = 98.6\,^\circ$F; $f(-10) = 14\,^\circ$F.

\textbf{3.} De $F = 1.8C + 32$: $C = \frac{F - 32}{1.8}$. Então
$50\,^\circ$F $\to \frac{18}{1.8} = 10\,^\circ$C, e $98.6\,^\circ$F
$\to \frac{66.6}{1.8} = 37\,^\circ$C.

\textbf{4.} $f(0) = 32 \neq 0$, e $f(20) = 68$ não é o dobro de
$f(10) = 50$: não é linear. O gráfico é uma reta que não passa pela
origem: $f$ é \emph{afim}, e não linear.

\textbf{5.} $1.8x + 32 = x$ dá $0.8x = -32$, então $x = -40$: a
quarenta abaixo de \omterm{def:g1:counting:zero}{zero}, os dois termômetros marcam $-40$. No gráfico,
a reta de $f$ cruza a reta diagonal $y = x$ exatamente ali.

\textbf{6.} $p_A(k) = 1.5k + 3$ e $p_B(k) = 2k$. Para $4$ km: $9$
euros contra $8$: B mais barato. Para $8$ km: $15$ contra $16$: A mais
barato.

\textbf{7.} $1.5k + 3 = 2k$ dá $k = 6$: em seis quilômetros, os dois
custam $12$ euros. Graficamente, a reta de B (pela origem, mais
inclinada) começa abaixo da de A e a cruza em $(6, 12)$: B ganha nas
corridas curtas, A ganha nas longas.

\textbf{8.} $N = 7 \times 20 - 30 = 110$ cricrilos por minuto. De
$82 = 7T - 30$: $T = \frac{112}{7} = 16\,^\circ$C.

\textbf{9.} $v(t) = 600 - 15t$. Valendo $240$: $600 - 15t = 240$,
$t = 24$ meses. Valendo $0$: $t = 40$ meses. Depois de $40$ meses, a
fórmula fica negativa --- mas o valor de um celular para em \omterm{def:g1:counting:zero}{zero}: todo
modelo afim tem um domínio em que faz sentido.

\textbf{10.} Compondo: $\times 1.25$ e depois $\times 0.8$ multiplica
por $1.25 \times 0.8 = 1$: o preço volta exatamente ao ponto de
partida. Nada de milagre: as porcentagens agem sobre referências
diferentes (o corte se aplica ao preço já aumentado) e, aqui, os
multiplicadores por acaso se cancelam --- a história geral, com a perda
embutida $1 - \frac{t^2}{10\,000}$, é o \cref{exo:g9:linfunc:10}.

\textbf{11.} \omterm{prop:g9:linfunc:affinegraph}{Coeficiente angular}: $\frac{17 - 20}{30} = -0.1$ cm por
minuto; $h(t) = 20 - 0.1t$. Ela se apaga em $h(t) = 0$: $t = 200$
minutos --- três horas e vinte minutos de luz.

\textbf{12.} O segundo partiu $10$ km à frente (coeficiente linear); o
primeiro é mais rápido ($24 > 18$ km/h). Alcance:
$24t = 10 + 18t$ dá $6t = 10$, $t = \frac53$ h $= 1$ h $40$ min, no
quilômetro $24 \times \frac53 = 40$.

\textbf{13.} $f(x) + g(x) = (2 - 0.5)x + (3 + 7) = 1.5x + 10$: afim,
com \omterm{prop:g9:linfunc:affinegraph}{coeficiente angular} igual à \emph{\omterm{def:g1:addition:def}{soma} dos \omterm{prop:g9:linfunc:affinegraph}{coeficientes angulares}}
e coeficiente linear igual à \omterm{def:g1:addition:def}{soma} dos coeficientes lineares --- fica
claro ao juntar os termos semelhantes em $(mx + p) + (m'x + p')$.

\textbf{14.} Se $m \neq 0$: exatamente uma solução, $x = -\frac pm$
--- uma reta não horizontal cruza o eixo uma vez. Se $m = 0$ e
$p \neq 0$: nenhuma solução --- uma reta horizontal acima ou abaixo do
eixo nunca o toca. Se $m = 0$ e $p = 0$: todo $x$ é solução --- a reta
\emph{é} o eixo.

\textbf{15.} Primeiro trio: taxas $\frac{9 - 5}{3 - 1} = 2$ e
$\frac{17 - 9}{7 - 3} = 2$: iguais, os pontos estão alinhados (sobre
$y = 2x + 3$). Segundo trio: $\frac{9-5}{2} = 2$, mas
$\frac{16 - 9}{3} = \frac73 \neq 2$: não estão alinhados. As \omterm{def:g9:linfunc:affine}{funções
afins} são exatamente aquelas de taxa de variação constante --- um único
número, o \omterm{prop:g9:linfunc:affinegraph}{coeficiente angular}, conta a história inteira; nas curvas, a
própria taxa muda de ponto para ponto, e persegui-la leva à derivada,
no volume do ensino médio.
\end{solution}
